В первой строке задано одно целое число $N$ $(2 \le N \le 3 \cdot 10^5)$ --- количество вершин дерева.

Во второй строке заданы через пробел $N$ целых чисел $p_1, p_2, \ldots, p_N$ $(0 \le p_i \le N)$ --- массив предков, где $p_i$ --- номер вершины-предка вершины $i$, при этом $p_i = 0$, если вершина $i$ является корнем дерева.

Гарантируется, что:

- ровно одно из чисел $p_i$ равно $0$;

- массив предков задаёт корректное дерево.

Вершины дерева пронумерованы от $1$ до $N$.
